\chapter*{Conclusion}
\addcontentsline{toc}{chapter}{Conclusion}
\markboth{\small\textit{Conclusion}}{\small\textit{Conclusion}}

La classification LC est à l'image d'un arbre dont les branches se sont développées selon des logiques propres au domaine du savoir qu'elles cherchent à couvrir.

Ce mémoire a observé le comportement du modèle de langage, confronté à cet arbre depuis le terrain particulier d'une bibliothèque patrimoniale spécialisée. À l'image de l'arbre lui-même, ses réactions se différencient branche par branche.

La question de ce mémoire appelle donc dès le départ une réponse qui ne peut pas être tout à fait affirmative ou tout à fait négative. Les résultats obtenus sur le corpus INHA (n~=~201) confirment cette intuition. Une simplification de la réponse est que la performance dépend de la classe considérée. Et cette hétérogénéité est un résultat observé dans nos tests.

\subsection*{Synthèse des résultats}

Deux effets résistent à une correction stricte statistique pour comparaisons multiples et peuvent être considérés comme statistiquement robustes au sens le plus strict retenu dans ce mémoire. Le premier est l'effet du plan de cotation réduit, curaté par la bibliothèque du musée d'Orsay, par rapport au barème officiel complet de la \textit{Library of Congress}~: en l'absence de \textit{few-shot}, le score moyen passe de 1,823 à 2,673 sur quatre ($p = 1,6\times10^{-5}$), soit l'écart le plus large observé dans l'ensemble de ce travail. Le second est l'effet de la priorisation des sujets RAMEAU sur le titre, sous plan réduit ($+0,318$, $p = 0,0035$).

Comme à l'image de l'arbre, ces deux effets robustes ne racontent cependant pas la même histoire une fois décomposés par grande classe. L'effet du plan de cotation, loin d'être uniforme, s'inverse pour la classe N~: cette classe, la plus peuplée et la plus difficile du corpus (34,5~\% d'exact-match), obtient de meilleurs résultats avec le barème \emph{complet} qu'avec le plan \emph{réduit}.

L'examen des erreurs révèle la cause de cette inversion~: sous plan réduit, 20 notices de la classe N sur 55 (36,4~\%) se voient attribuer exactement la même cote (\texttt{N7560-8266}), quel que soit leur sujet réel, parce que le plan ne propose qu'une entrée unique pour toute cette portion du barème~; ce repli disparaît entièrement dès lors que le modèle dispose du texte intégral, plus richement subdivisé pour cette même zone.

Deux branches du même arbre, soumises au même levier expérimental, ont donc évolué en sens opposés~: la réduction du plan a affiné le comportement du modèle sur les branches NK et AM (arts décoratifs, musées), tandis qu'elle l'a poussé, sur la branche N, vers un repli homogène et peu informatif. La curation experte, loin de toujours faciliter la tâche du modèle, peut ainsi produire l'effet inverse de celui recherché lorsqu'elle laisse une zone insuffisamment différenciée.

Un second résultat, non anticipé au moment de la conception du protocole expérimental, mérite d'être souligné~: l'ajout d'exemples \textit{few-shot} au \textit{prompt} dégrade les résultats sous plan réduit ($-0,279$, $p = 0,027$), la meilleure configuration mesurée dans ce mémoire étant, toutes conditions confondues, le plan réduit \emph{sans} few-shot (66,2~\% d'exact-match). L'élargissement du pool d'exemples produit lui aussi un effet contrasté selon les branches~: bénéfique en moyenne sous plan complet, mais dégradant pour la classe NX. Aucun enrichissement du \textit{prompt}, quel qu'il soit, ne doit être présumé bénéfique sans être mesuré explicitement, et sur chaque branche concernée plutôt que sur la seule moyenne globale.

\subsection*{Limites}

Trois limites structurent la portée de ces résultats et méritent d'être nommées clairement, au-delà de celles déjà détaillées.

La première concerne la source du barème utilisé pour mesurer la proximité des erreurs. Ce mémoire a d'abord mobilisé un résumé du barème officiel de la \textit{Library of Congress}, largement diffusé et réutilisé dans la littérature technique sur le sujet (\textit{outline}, quelques centaines d'entrées), avant de découvrir, en cours de rédaction, qu'il sous-représentait considérablement certaines classes~: la classe AM, par exemple, n'y compte que quatre entrées, contre près d'une cinquantaine dans le texte intégral des \textit{schedules} officiels. Cette découverte a conduit à extraire directement le texte intégral pour les classes couvrant l'essentiel du corpus (AM et la famille N, soit 92~\% des notices évaluées), et à revoir en conséquence plusieurs interprétations initiales de ce mémoire.

La seconde limite, qui découle directement de la première, touche au score de proximité lui-même. Une fois le texte intégral substitué au résumé pour le contenu montré au modèle, la hiérarchie utilisée pour \emph{mesurer} la proximité des erreurs n'a pu être que partiellement reconstituée~: les catégories intermédiaires du barème officiel (les regroupements continentaux au sein d'une subdivision par pays, par exemple) n'ont donc pas pu être restituées par la seule reconstruction automatique tentée dans ce mémoire. Deux sous-classes voisines peuvent ainsi se retrouver, pour le calcul du score, aussi éloignées l'une de l'autre que deux sous-classes sans rapport. Ce biais ne remet pas en cause les différences mesurées entre elles, mais limite la lecture du score absolu de chaque condition prise isolément et rend la mesure du score moins précise.

La troisième limite tient à l'absence de validation humaine systématique des scores produits. Le score de proximité automatique, reste un indicateur indirect de la justesse réelle d'une proposition de cotation~; seul un bibliothécaire spécialisé peut trancher, sur les cas ambigus, si un score intermédiaire ou nul reflète une vraie erreur sémantique ou un artefact de mesure. 

\subsection*{Apports et perspectives}

Au-delà de ces limites, ce mémoire ouvre plusieurs pistes de travail directement actionnables.

La première consisterait à étendre l'extraction du texte intégral du barème officiel aux classes du périmètre BM'O non encore couvertes (BH, C, CC, CZ, D, DA, GN, NB, NC, NE, SB, TR), afin de disposer d'une mesure homogène sur l'ensemble du corpus plutôt que sur les seules classes majoritaires. La reconstruction d'une hiérarchie plus fidèle, en particulier des niveaux intermédiaires actuellement absents, en serait un prolongement naturel~: elle supposerait de repartir du texte brut des \textit{schedules} en conservant l'indentation visuelle d'origine, plutôt que de reconstruire la hiérarchie par le seul confinement numérique des plages, comme cela a été tenté dans le cadre de ce mémoire sans résultat concluant. Une métrique alternative, fondée sur l'écart numérique normalisé entre deux cotes plutôt que sur une hiérarchie explicite, constitue une autre piste~: la LCC organisant ses sous-classes de façon globalement continue au sein d'une même grande classe, la proximité numérique pourrait s'y révéler un indicateur de proximité thématique suffisant, sans dépendre d'un barème hiérarchique complet difficile à obtenir au niveau de détail requis.

La seconde piste porte sur la généralisation au-delà du modèle et du corpus testés. L'effet d'attracteur documenté sur la classe N, bien que sensible au plan de cotation utilisé, reste un comportement propre au modèle Mistral évalué ici~; sa réplication sur d'autres modèles de langage permettrait de déterminer s'il s'agit d'un biais transversal aux LLM en tâche de classification fine, ou d'une singularité de ce modèle précis. De même, l'évaluation menée sur le corpus INHA constitue un proxy, choisi en l'absence de vérité terrain disponible sur le fonds réel de la BM'O~; une évaluation ultérieure sur un échantillon du fonds réel, une fois celui-ci constitué, permettrait de vérifier si les effets mesurés ici se maintiennent à cette échelle et sur des notices dont les métadonnées suivent les conventions propres à la bibliothèque plutôt que celles d'un fonds de substitution.

Enfin, le résultat le plus inattendu de ce mémoire est l'effet dégradant du \textit{few-shot} sous plan réduit. Cet effet mériterait une investigation plus poussée, qui dépasse le cadre de cette étude~: identifier si cet effet tient à la nature des exemples fournis, à leur nombre, ou à une interférence plus générale entre exemples et instructions explicites dans le \textit{prompt}, permettrait de formuler des recommandations plus fines que le simple retrait de ce levier dans la configuration recommandée en~\S\ref{sec:workflow}.

Ce que ce mémoire retient, en dernière instance, n'est donc pas qu'un périmètre documentaire réduit facilite ou complique, en bloc, l'automatisation de la cotation, mais que la réduction opérée préserve une granularité suffisante pour guider le modèle plutôt que de le pousser vers un repli par défaut.

Darwin observait que l'arbre du vivant ne se ramifie pas selon un plan préconçu, mais selon la pression locale que chaque branche rencontre~; ce mémoire suggère qu'il en va de même pour l'arbre de la classification, une fois confié à un modèle de langage. Aucun réglage unique ne vaut pour l'ensemble de ses branches, et une bibliothèque patrimoniale spécialisée qui souhaiterait appliquer une démarche comparable devrait non seulement curer son périmètre documentaire, mais vérifier, branche par branche, que cette curation ne produit pas l'effet inverse de celui recherché.